%% ============================================================
%%  projects.tex — Applied Mathematics Project Portfolio
%%  Fonts: Palatino  |  No fontawesome5 / lmodern dependency
%%  KEY FIX: \tag renamed to \cvtag (amsmath owns \tag)
%%  Compile: pdflatex projects.tex  (run twice)
%% ============================================================

\documentclass[10pt,letterpaper]{article}

%% ── Geometry ────────────────────────────────────────────────
\usepackage[top=0.65in,bottom=0.65in,left=0.70in,right=0.70in]{geometry}

%% ── Encoding & Font ─────────────────────────────────────────
\usepackage[T1]{fontenc}
\usepackage[utf8]{inputenc}
\usepackage{palatino}
\usepackage{microtype}

%% ── Math ────────────────────────────────────────────────────
\usepackage{amsmath,amssymb,amsthm}

%% ── Colors ──────────────────────────────────────────────────
\usepackage[dvipsnames,svgnames]{xcolor}
\definecolor{NavyDeep}{HTML}{1B2A4A}
\definecolor{AccentTeal}{HTML}{2A6E7C}
\definecolor{RuleGray}{HTML}{C8CDD5}
\definecolor{SubtleGray}{HTML}{6B7280}
\definecolor{TagBg}{HTML}{E8F4F6}

%% ── Hyperlinks ──────────────────────────────────────────────
\usepackage[hidelinks]{hyperref}
\hypersetup{colorlinks=true,urlcolor=AccentTeal,linkcolor=NavyDeep}

%% ── Lists ───────────────────────────────────────────────────
\usepackage{enumitem}
\setlist[itemize]{
  leftmargin=1.4em,itemsep=1.5pt,parsep=0pt,topsep=2pt,
  label={\color{AccentTeal}\textbullet}
}

%% ── Tables ──────────────────────────────────────────────────
\usepackage{booktabs,tabularx,array}

%% ── Colored boxes ───────────────────────────────────────────
\usepackage{tcolorbox}
\tcbuselibrary{skins,breakable}

%% ── Icons ───────────────────────────────────────────────────
\usepackage{marvosym}   % \Email etc.
\usepackage{pifont}     % \ding{} checkmarks etc.

%% ── Page-break control ──────────────────────────────────────
\usepackage{needspace}

%% ── Misc ────────────────────────────────────────────────────
\usepackage{parskip,etoolbox}
\setlength{\parskip}{0pt}
\pagestyle{empty}

%% ============================================================
%%  CUSTOM COMMANDS
%% ============================================================

\newcommand{\pfsection}[1]{%
  \needspace{6\baselineskip}%
  \vspace{6pt}%
  {\color{NavyDeep}\rule{\linewidth}{0.8pt}\par}%
  \vspace{-2pt}%
  {\color{NavyDeep}\large\bfseries\MakeUppercase{#1}\par}%
  \vspace{3pt}%
}

%% Project card header: \projecthead{Title}{Subtitle}{Date}
\newcommand{\projecthead}[3]{%
  \noindent
  \begin{tabularx}{\linewidth}{@{}X r@{}}
    {\large\bfseries\color{NavyDeep}#1} & {\small\color{SubtleGray}#3}\\[-2pt]
    {\small\itshape\color{AccentTeal}#2} & \\
  \end{tabularx}%
  \vspace{2pt}%
}

%% Method/tool tag  — NOTE: named \cvtag to avoid clash with amsmath \tag
\newcommand{\cvtag}[1]{%
  \mbox{\small\bfseries\color{AccentTeal}\texttt{[#1]}}\,%
}

%% Key result callout box
\newcommand{\keyresult}[1]{%
  \vspace{3pt}%
  \noindent\colorbox{TagBg}{%
    \parbox{\dimexpr\linewidth-2\fboxsep\relax}{%
      \small\color{NavyDeep}%
      \textbf{Key result:}\enspace#1%
    }%
  }%
  \vspace{3pt}%
}

%% Inline math shortcut
\newcommand{\eq}[1]{$#1$}

%% ============================================================
%%  DOCUMENT
%% ============================================================
\begin{document}

%% ── Page header ─────────────────────────────────────────────
\begin{center}
  {\fontsize{20}{24}\selectfont\bfseries\color{NavyDeep} Project Portfolio}\\[3pt]
  {\color{AccentTeal}\rule{3cm}{1pt}}\\[4pt]
  {\small\color{SubtleGray}
    Your Full Name\enspace\textbar\enspace
    \href{mailto:you@university.edu}{you@university.edu}\enspace\textbar\enspace
    \href{https://github.com/yourhandle}{github.com/yourhandle}\enspace\textbar\enspace
    \href{https://yoursite.com}{yoursite.com}
  }
\end{center}
\vspace{4pt}

%% ── Project 1 ───────────────────────────────────────────────
\pfsection{Selected Research Projects}

\projecthead
  {Adaptive Multilevel Monte Carlo for Stochastic PDEs}
  {Ph.D.\ dissertation project, University of Somewhere}
  {2021--Present}

\noindent\cvtag{Numerical Analysis} \cvtag{UQ} \cvtag{C++/MPI} \cvtag{SIAM J.\ Sci.\ Comput.}

\vspace{4pt}
\textbf{Problem.}\enspace
UQ for subsurface flow requires statistics of solutions to elliptic PDEs
\eq{-\nabla\cdot(a(\mathbf{x},\omega)\nabla u)=f} with log-normal random coefficient \eq{a}.
Standard Monte Carlo costs \eq{\mathcal{O}(\varepsilon^{-2}h^{-d})} for RMSE \eq{\varepsilon}.

\textbf{Contribution.}\enspace
I designed an \emph{adaptive} multilevel estimator that selects mesh levels
\eq{\{h_\ell\}} and sample counts \eq{\{N_\ell\}} online by minimising cost subject to
a variance budget, using a Gaussian process surrogate to predict level contributions.

\begin{itemize}
  \item Proved optimal complexity \eq{\mathcal{O}(\varepsilon^{-2}|\log\varepsilon|^2)}
        without \emph{a priori} knowledge of coefficient regularity.
  \item Implemented in C++17/MPI with FEniCS and PETSc; near-linear weak scaling to
        512 MPI tasks on Stampede2.
  \item \eq{3\times} wall-clock speedup over standard MLMC on a 3-D Darcy flow benchmark
        with Mat\'{e}rn-\eq{3/2} coefficient.
\end{itemize}

\keyresult{Paper accepted, \textit{SIAM J.\ Sci.\ Comput.}\ 2024. Code: \href{https://github.com/yourhandle/adaptive-mlmc}{\texttt{github.com/yourhandle/adaptive-mlmc}}.}

\vspace{8pt}

%% ── Project 2 ───────────────────────────────────────────────
\projecthead
  {Sharp Convergence Analysis of Nonconvex SGD}
  {Research internship, National Lab}
  {Summer 2022}

\noindent\cvtag{Optimization Theory} \cvtag{ML Theory} \cvtag{Python/JAX} \cvtag{arXiv preprint}

\vspace{4pt}
\textbf{Problem.}\enspace
Convergence guarantees for SGD on nonconvex objectives remain loose;
prior bounds require bounded gradients or structural assumptions rarely satisfied in practice.

\textbf{Contribution.}\enspace
I established a sharp \eq{O(k^{-1/2})} rate for
\eq{\mathbb{E}[\|\nabla f(x_k)\|^2]} under \eq{(L,\sigma)}-smoothness — strictly
weaker than standard \eq{L}-smoothness — via a Lyapunov function accounting for
batch-size-dependent variance.

\begin{itemize}
  \item Derived tight lower bounds confirming \eq{O(k^{-1/2})} is unimprovable without
        additional curvature information.
  \item Validated on CIFAR-10/100; empirical convergence matched predicted rates across
        6 architectures.
  \item Compared AdaGrad, Adam, SGD-M; identified regimes where adaptive methods
        provably dominate.
\end{itemize}

\keyresult{Preprint arXiv:2401.00000; under revision for NeurIPS 2024.}

\vspace{8pt}

%% ── Project 3 ───────────────────────────────────────────────
\projecthead
  {Sparse-Grid Collocation for High-Dimensional Bayesian Inversion}
  {Independent study / course project}
  {Spring 2023}

\noindent\cvtag{Bayesian Inference} \cvtag{Sparse Grids} \cvtag{Python} \cvtag{SciPy/Chaospy}

\vspace{4pt}
\textbf{Problem.}\enspace
Bayesian inversion of distributed parameter fields leads to posteriors on
high-dimensional (\eq{d=50\text{--}200}) spaces; MCMC is prohibitive when each
forward solve costs \eq{\sim}1 CPU-minute.

\textbf{Approach.}\enspace
I built a surrogate-accelerated pipeline: a Smolyak sparse-grid interpolant
approximates the log-likelihood offline (\eq{\sim 10^3} PDE solves), then
Metropolis-Hastings samples the surrogate at negligible cost.

\begin{itemize}
  \item \eq{<}5\% relative error in posterior mean and variance for a 50-D
        groundwater conductivity inversion.
  \item Clenshaw--Curtis nodes reduce interpolation error by \eq{2\times} over
        Legendre nodes for smooth likelihoods.
  \item Reusable Python module with FEniCS interface; tested on two published benchmarks.
\end{itemize}

\vspace{8pt}

%% ── Additional Projects ─────────────────────────────────────
\pfsection{Additional Projects}

\projecthead
  {Finite-Element Solver for Incompressible Navier--Stokes}
  {Senior honors thesis, State University}
  {2019--2020}

\noindent\cvtag{FEM} \cvtag{Fluid Dynamics} \cvtag{C++} \cvtag{deal.II}

\vspace{4pt}
Taylor--Hood \eq{(\mathbf{P}_2/P_1)} mixed FEM for 2-D incompressible flow using deal.II;
second-order velocity convergence verified via manufactured solutions.
Lid-driven cavity flow at \eq{\mathrm{Re}=400,\,1000,\,3200} matches Ghia et al.\
(1982) to 4 significant figures.

\vspace{6pt}

\projecthead
  {Computer-Assisted Hopf Bifurcation Proof}
  {REU project, Another University}
  {Summer 2019}

\noindent\cvtag{Dynamical Systems} \cvtag{Interval Arithmetic} \cvtag{MATLAB/Intlab}

\vspace{4pt}
Used INTLAB to rigorously enclose Jacobian eigenvalues of a 3-species Lotka--Volterra
system and certify a sign change in the real part, constituting a computer-assisted
proof of Hopf bifurcation at \eq{\mu^*=0.7134\pm10^{-4}}.

\vspace{6pt}

%% ── Skills summary ──────────────────────────────────────────
\pfsection{Technical Stack}

\noindent{\small
\begin{tabularx}{\linewidth}{@{}l X@{}}
\toprule
\textbf{Domain} & \textbf{Tools \& Methods}\\
\midrule
Numerical PDE   & FEniCS, deal.II, PETSc, Firedrake, FDM, FEM, Spectral methods\\
HPC             & C++17, MPI, OpenMP, SLURM, vectorised NumPy, JAX JIT\\
UQ / Stoch.\    & MLMC, Sparse Grids, Polynomial Chaos, MCMC, Kalman filters\\
Optimisation    & SciPy, JAX (autodiff), L-BFGS, proximal methods, ADMM\\
ML              & PyTorch, scikit-learn, Gaussian processes (GPyTorch)\\
Languages       & Python, C/C++, MATLAB, Julia, R, Bash\\
\bottomrule
\end{tabularx}
}

\end{document}
